% AA-45 | Kryteria hiperwykładnicze i całki nieelementarne.
% Źródła: AA-43, AA-44, sources/pl/galois(2).tex.
% Stan: nowe opracowanie, przed niezależnymi audytami i finalną kontrolą PDF.
\chapter{Kryteria hiperwykładnicze i całki nieelementarne}
\chaptermark{Całki nieelementarne}
\label{aa:ch:nonelementary-integrals}

W rozdziale~\ref{aa:ch:liouville-elementary-integration}
udowodniliśmy twierdzenie Liouville'a: element $f$ ciała
różniczkowego ma elementarną funkcję pierwotną bez nowych stałych
wtedy i tylko wtedy, gdy jest sumą pochodnej elementu tego ciała
i skończonej kombinacji pochodnych logarytmicznych
o stałych współczynnikach. Daje to kryterium strukturalne,
które trzeba przekształcić w równania poddające się badaniu.
Najpierw zrobimy to dla rozszerzeń wykładniczych,
a potem przeanalizujemy cztery różne całki nieelementarne.

\section{Funkcje elementarne i ich lokalna interpretacja}

Przy funkcjach zespolonych używamy zarodków funkcji
meromorficznych w punkcie, w którym skończona liczba
rozważanych wyrażeń jest określona; dla pierwiastków
i logarytmów wybieramy lokalne gałęzie. Ciało takich
zarodków zawierające wszystkie stałe zespolone ma
dokładnie $\mathbb C$ jako ciało stałych:
zarodek o zerowej pochodnej jest stały.
Dlatego dowód nieelementarności w rozszerzeniu
bez nowych stałych dotyczy również zwykłych
lokalnych wzorów elementarnych.

Z rzeczywistego wzoru elementarnego na otwartym
przedziale można w zwykłym punkcie otrzymać
lokalny wzór zespolony, wybierając holomorficzne
gałęzie jego składników. Taką interpretację
stosujemy poniżej: wykluczenie lokalnej
zespolonej funkcji pierwotnej elementarnej
wyklucza wzór na rzeczywistym przedziale.
Nie twierdzimy, że logarytmy i pierwiastki
są globalnie jednowartościowe.

\section{Elementy hiperwykładnicze}

\begin{definicja}[Element hiperwykładniczy]
\label{aa:def:hyperexponential}
Niech $(F,D)$ będzie ciałem różniczkowym
charakterystyki $0$. Niezerowy element $t$
rozszerzenia różniczkowego $L/F$ nazywamy
\emph{hiperwykładniczym nad $F$}, jeśli
$Dt/t=h$ dla pewnego $h\in F$.
W kryteriach poniżej zakładamy dodatkowo,
że $t$ jest przestępny nad $F$ i
$C_{F(t)}=C_F$. Te dwa założenia są
istotne, a nie są częścią samego równania $Dt=ht$.
\end{definicja}

\begin{lemat}[Wykładnicze wielomianów]
\label{aa:lem:exp-polynomial-constant-field}
Niech $p\in\mathbb C[x]$ będzie niestałym
wielomianem i niech $K_0=\mathbb C(x)$
z $Dx=1$. Wówczas element $t=e^{p(x)}$,
rozumiany jako zarodek analityczny,
jest przestępny nad $K_0$, spełnia
$Dt=p'(x)t$, a $C_{K_0(t)}=\mathbb C$.
\end{lemat}
\begin{proof}
Każdy $u\in\mathbb C(x)^\times$ ma,
po rozkładzie licznika i mianownika,
postać $u=c\prod_{\nu=1}^{N}(x-a_\nu)^{m_\nu}$,
gdzie $c\in\mathbb C^\times$, $a_\nu\in\mathbb C$
są parami różne, a $m_\nu\in\mathbb Z$.
Stąd
\[
 \frac{Du}{u}=\sum_{\nu=1}^N\frac{m_\nu}{x-a_\nu}.
\]
Ta pochodna logarytmiczna jest właściwą
funkcją wymierną: jej część wielomianowa
w $x$ jest zerowa. Nie może więc być
niezerowym wielomianem.

Przypuśćmy, że $t$ jest algebraiczny nad
$K_0$ i niech $d=[K_0(t):K_0]$.
Stosując tożsamość dla normy z
Twierdzenia~\ref{aa:tw:differential-trace-norm}
do elementu $t$, otrzymujemy
\[
 \frac{D N_{K_0(t)/K_0}(t)}
        {N_{K_0(t)/K_0}(t)}
 =\operatorname{Tr}_{K_0(t)/K_0}(Dt/t)
 =d\,p'(x).
\]
Po lewej jest pochodna logarytmiczna
niezerowej funkcji wymiernej,
a po prawej niezerowy wielomian:
$p$ jest niestały i charakterystyka wynosi $0$.
Otrzymaliśmy sprzeczność.

Pozostaje zbadać stałe. Na mocy
Lematu~\ref{aa:lem:exponential-extension-constants}
nowa stała w $\mathbb C(x,t)$ istnieje
dokładnie wtedy, gdy
$Du/u=-m p'(x)$ dla pewnego
$u\in\mathbb C(x)^\times$ i
$m\in\mathbb Z\setminus\{0\}$.
Prawa strona jest niezerowym wielomianem,
a lewa właściwą funkcją wymierną.
Takiej pary nie ma, czyli
$C_{\mathbb C(x,t)}=\mathbb C$.
\end{proof}

\begin{uwaga}[Rola założeń]
\label{aa:uw:hyperexponential-obstructions}
Równanie $Dt=ht$ nie gwarantuje samo
przez się zachowania stałych.
Jeśli $K_0=\mathbb C(x)$, $h=1/x$
i $t$ jest przestępny z $Dt=t/x$,
to $D(t/x)=0$, a $t/x\notin K_0$.
Również przy $h=0$ przestępny generator
mający $Dt=0$ jest nową stałą.
Z kolei $e^{p(x)}$ z $p$ stałym
już należy do $\mathbb C$.
\end{uwaga}

\section{Pełne kryterium dla skończonych sum wykładniczych}

\begin{twierdzenie}[Kryterium hiperwykładnicze]
\label{aa:tw:hyperexponential-criterion}
Niech $h\in\mathbb C(x)$, a $t$ będzie
przestępny nad $\mathbb C(x)$ i spełnia
$Dt=ht$ oraz $C_{\mathbb C(x,t)}=\mathbb C$.
Dla $s\in\mathbb C(x)$ następujące warunki
są równoważne:
\begin{enumerate}
\item element $st$ ma elementarną funkcję
pierwotną nad $\mathbb C(x,t)$, bez nowych stałych;
\item istnieje $r\in\mathbb C(x)$ takie, że
\begin{equation}\label{aa:eq:hyperexp-ode}
 r'+hr=s.
\end{equation}
\end{enumerate}
Jeżeli te warunki zachodzą, element $rt$
jest funkcją pierwotną $st$.
\end{twierdzenie}
\begin{proof}
Jeżeli $r'+hr=s$, to
$D(rt)=(r'+hr)t=st$,
co daje elementarną funkcję pierwotną
w samym $\mathbb C(x,t)$.

Załóżmy odwrotnie istnienie elementarnej
funkcji pierwotnej. Z twierdzenia
\ref{aa:tw:liouville-elementary-integration}
dla ciała $\mathbb C(x,t)$ otrzymujemy
\[
 st=Dv+\sum_{j=1}^N c_j\frac{Du_j}{u_j},
 \qquad v\in\mathbb C(x,t),\quad
 u_j\in\mathbb C(x,t)^\times,\quad
 c_j\in\mathbb C.
\]
Niech $P=\operatorname{pp}_t(v)\in\mathbb C(x)[t]$.
Lemat~\ref{aa:lem:liouville-polynomial-parts}
daje po zastosowaniu $\operatorname{pp}_t$:
$st=DP+b(x)$, gdzie $b\in\mathbb C(x)$;
wszystkie pochodne logarytmiczne wnoszą
jedynie wyrazy niezależne od $t$.
Oznaczmy przez $r$ współczynnik przy
$t$ w $P$, przyjmując $r=0$, jeśli go nie ma.
Współczynnik przy $t$ w $DP$ wynosi
$r'+hr$. Przestępność $t$ zapewnia
jednoznaczność porównania współczynników,
więc otrzymujemy~\eqref{aa:eq:hyperexp-ode}.
\end{proof}

\begin{twierdzenie}[Kryterium dla wielomianów Laurenta]
\label{aa:tw:laurent-hyperexp-criterion}
Przy założeniach
Twierdzenia~\ref{aa:tw:hyperexponential-criterion},
niech $S\subset\mathbb Z$ będzie zbiorem
skończonym, a dla $k\in S$ niech
$s_k\in\mathbb C(x)$. Element
$f=\sum_{k\in S}s_k(x)t^k\in\mathbb C(x,t)$
ma elementarną funkcję pierwotną bez nowych
stałych wtedy i tylko wtedy, gdy dla
każdego $k\in S\setminus\{0\}$ istnieje
$r_k\in\mathbb C(x)$ spełniające
\begin{equation}\label{aa:eq:laurent-ode}
 r_k'+khr_k=s_k.
\end{equation}
Jeśli $0\notin S$, przyjmujemy $s_0=0$.
Współczynnik $s_0$ nie nakłada dodatkowego
warunku: każda funkcja wymierna nad
$\mathbb C$ ma elementarną funkcję pierwotną.
\end{twierdzenie}
\begin{proof}
Jeśli istnieją wszystkie $r_k$,
to $D(r_kt^k)=(r_k'+khr_k)t^k$.
Pozostałą funkcję wymierną $s_0$ całkujemy
elementarnie: po rozkładzie na ułamki proste
nad $\mathbb C$ składniki wielomianowe
i bieguny rzędu co najmniej $2$ mają
wymierne funkcje pierwotne, a bieguny proste
dają stałe razy logarytmy. Te ostatnie
można dołączać bez nowych stałych zgodnie
z dowodem dostateczności twierdzenia
\ref{aa:tw:liouville-elementary-integration}.
To dowodzi wystarczalności.

Dla konieczności stosujemy twierdzenie
Liouville'a do $f$. Dla ustalonego $k>0$
porównujemy współczynniki przy $t^k$
w częściach wielomianowych po obu stronach
przedstawienia Liouville'a. Na mocy
Lematu~\ref{aa:lem:liouville-polynomial-parts}
pochodne logarytmiczne nie mają
dodatnich potęg $t$, dlatego
współczynnik $r_k\in\mathbb C(x)$ przy
$t^k$ w $\operatorname{pp}_t(v)$
spełnia~\eqref{aa:eq:laurent-ode}.
Współczynnik może być zerowy, jeśli
brak odpowiedniego wyrazu.

Dla ustalonego $k<0$ wprowadzamy
$w=t^{-1}$. Wtedy
$\mathbb C(x,w)=\mathbb C(x,t)$,
$Dw=-hw$ i $w$ jest przestępny.
Zastosowanie poprzedniego argumentu
do współczynnika przy $w^{-k}$ daje
$r_k'+(-k)(-h)r_k=s_k$, czyli
\eqref{aa:eq:laurent-ode}.
Nie twierdzimy, że pochodne logarytmiczne
nie mają ujemnych potęg w rozwinięciu
względem $t^{-1}$; dlatego osobna zmiana
generatora jest konieczna.
\end{proof}

\begin{przyklad}[Wzory elementarne i ich test]
\label{aa:prz:hyperexp-positive-controls}
Dla $t=e^{-x^2}$ mamy $h=-2x$.
Równanie $r'-2xr=x$ ma rozwiązanie
$r=-1/2$, a więc
$\int x e^{-x^2}\,\mathrm dx=-\frac12e^{-x^2}+C$.
Dla $t=e^{x^2}$ równanie
$r'+2xr=x$ ma rozwiązanie $r=1/2$.
Również
\[
 D\left(\frac{e^x}{x}\right)
 =e^x\left(\frac1x-\frac1{x^2}\right)
\]
jest przykładem dodatnim mimo biegunów
współczynników. Samo występowanie
wykładniczej nie przesądza zatem
o nieelementarności funkcji pierwotnej.
\end{przyklad}

\section{Przeszkody wymierne: bieguny i stopnie}

\begin{lemat}[Przeszkoda bieguna prostego]
\label{aa:lem:simple-pole-obstruction}
Niech $h,s\in\mathbb C(x)$ i niech
$a\in\mathbb C$. Załóżmy, że $h$
nie ma bieguna w $a$, natomiast $s$
ma w $a$ biegun dokładnie rzędu $1$.
Wtedy równanie $r'+hr=s$
nie ma rozwiązania $r\in\mathbb C(x)$.
\end{lemat}
\begin{proof}
Jeśli $r$ jest regularna w $a$,
to $r'+hr$ też jest regularna,
wbrew założeniu o $s$.
Jeśli $r$ ma biegun rzędu $m\ge1$,
to w lokalnej współrzędnej $z=x-a$
jej rozwinięcie zaczyna się od
$c z^{-m}$, $c\ne0$. Pochodna
$r'$ ma biegun rzędu $m+1$ z niezerowym
współczynnikiem $-mc$, gdyż
charakterystyka wynosi $0$.
Ponieważ $h$ jest regularna, $hr$
ma biegun rzędu co najwyżej $m$.
Nie może więc skasować wyrazu
$-mc z^{-m-1}$. Lewa strona ma
biegun rzędu $m+1\ge2$, sprzeczność.
\end{proof}

\begin{lemat}[Przeszkoda stopnia dla wielomianów]
\label{aa:lem:polynomial-degree-obstruction}
Niech $h,s\in\mathbb C[x]$,
$h\ne0$ i $\deg h\ge1$. Jeśli
$r\in\mathbb C(x)$ spełnia
$r'+hr=s$, to $r\in\mathbb C[x]$.
Jeżeli $s\ne0$ i $\deg s<\deg h$,
równanie nie ma rozwiązania wymiernego.
\end{lemat}
\begin{proof}
Gdyby $r$ miała skończony biegun
rzędu $m\ge1$, pochodna $r'$
miałaby biegun rzędu $m+1$,
a $hr$ co najwyżej $m$,
bo $h$ jest wielomianem.
Nie mogłyby się skasować,
tymczasem $s$ nie ma biegunów.
Zatem $r$ nie ma biegunów skończonych
i jako funkcja wymierna jest wielomianem.

Jeżeli $r$ jest niezerowym wielomianem
stopnia $d\ge0$, to $hr$ ma stopień
$d+\deg h$, podczas gdy $r'$
ma stopień co najwyżej $d-1$
(lub jest zerem). Stąd
$\deg(r'+hr)=d+\deg h\ge\deg h$.
Dla $\deg s<\deg h$ otrzymujemy
sprzeczność. Element $r=0$ nie
rozwiązuje równania z $s\ne0$.
\end{proof}

\begin{uwaga}[Założenie regularności $h$]
\label{aa:uw:pole-obstruction-edge}
W Lemacie~\ref{aa:lem:simple-pole-obstruction}
nie można opuścić założenia o regularności $h$.
Dla $r=1/x$ i $h=1/x$ zachodzi
$r'+hr=-1/x^2+1/x^2=0$:
dwa bieguny rzędu $2$ mogą się skasować
wtedy, gdy $h$ sama ma biegun w $0$.
\end{uwaga}

\section{Całka Gaussa}

\begin{twierdzenie}[Nieelementarność całki Gaussa]
\label{aa:tw:gaussian-nonelementary}
Funkcja $e^{-x^2}$ nie ma elementarnej
funkcji pierwotnej na żadnym niepustym
otwartym przedziale rzeczywistym.
Algebraicznie element $t$ z $Dt=-2xt$
nie ma elementarnej funkcji pierwotnej
nad $\mathbb C(x,t)$ bez nowych stałych.
\end{twierdzenie}
\begin{proof}
Lemat~\ref{aa:lem:exp-polynomial-constant-field}
dla $p=-x^2$ uzasadnia przestępność
$t=e^{-x^2}$ oraz brak nowych stałych.
Twierdzenie~\ref{aa:tw:hyperexponential-criterion}
z $s=1$ i $h=-2x$ wymagałoby
rozwiązania $r'-2xr=1$
w $\mathbb C(x)$. Lemat
\ref{aa:lem:polynomial-degree-obstruction}
wyklucza takie rozwiązanie, gdyż
$\deg 1=0<1=\deg(-2x)$.
Nieelementarność lokalnego wzoru
zespolonego wyklucza też wzór
na dowolnym przedziale rzeczywistym.
\end{proof}

\begin{przyklad}[Funkcja błędu i całka oznaczona]
\label{aa:prz:gaussian-erf-value}
Brak elementarnej funkcji pierwotnej nie oznacza,
że funkcja pierwotna nie istnieje.
Można zdefiniować funkcję
$\operatorname{erf}(x)=
(2/\sqrt{\pi})\int_0^x e^{-u^2}\,\mathrm du$.
Jest ona różniczkowalna, a jej pochodna
wynosi $(2/\sqrt{\pi})e^{-x^2}$.

Ponadto całka niewłaściwa
$I=\int_0^\infty e^{-x^2}\,\mathrm dx$
jest zbieżna: dla $x\ge1$ zachodzi
$e^{-x^2}\le e^{-x}$. Tonelli pozwala
pomnożyć dodatnie całki, a zmiana
zmiennych na biegunowe w pierwszej
ćwiartce płaszczyzny daje
\[
 I^2=\int_0^\infty\int_0^\infty
 e^{-(x^2+y^2)}\,\mathrm dx\,\mathrm dy
 =\int_0^{\pi/2}\int_0^\infty
 e^{-\rho^2}\rho\,\mathrm d\rho\,\mathrm d\theta
 =\frac{\pi}{4}.
\]
Ponieważ $I>0$, otrzymujemy
$I=\sqrt{\pi}/2$. Wartość oznaczona
może więc być elementarna mimo
nieelementarności funkcji pierwotnej.
\end{przyklad}

\section{Całka sinusowa}

\begin{twierdzenie}[Nieelementarność całki sinusowej]
\label{aa:tw:sine-integral-nonelementary}
Funkcja $\sin x/x$, z wartością $1$
w punkcie $x=0$, nie ma elementarnej
funkcji pierwotnej na żadnym
niepustym otwartym przedziale rzeczywistym.
\end{twierdzenie}
\begin{proof}
Przyjmijmy $t=e^{ix}$, więc
$Dt=it$. Lemat
\ref{aa:lem:exp-polynomial-constant-field}
z $p(x)=ix$ zapewnia przestępność
i brak nowych stałych.
Korzystamy z pojedynczego generatora:
\[
 \frac{\sin x}{x}
 =\frac{t-t^{-1}}{2ix}.
\]
Nie zakładamy niezależności
$t$ i $t^{-1}$; ich iloczyn wynosi $1$.
Twierdzenie~\ref{aa:tw:laurent-hyperexp-criterion}
dla współczynnika przy dodatniej
potędze $t$ wymuszałoby
$r'+ir=1/(2ix)$ dla
$r\in\mathbb C(x)$.
Po położeniu $R=2ir$
dostalibyśmy $R'+iR=1/x$.
Jest to niemożliwe na mocy
Lematu~\ref{aa:lem:simple-pole-obstruction}
w punkcie $0$ z $h=i$.
Nawet jeśli sam iloraz $\sin x/x$
jest regularny w zerze, wymuszone
równanie wymierne ma tam biegun prosty.
Sprzeczność dowodzi tezy.
\end{proof}

\begin{przyklad}[Całka Dirichleta i funkcja $\operatorname{Si}$]
\label{aa:prz:sine-integral-dirichlet}
Funkcję sinusową definiuje się przez
$\operatorname{Si}(x)=
\int_0^x (\sin u/u)\,\mathrm du$,
przyjmując wartość $1$ funkcji podcałkowej
w zerze. Wyznaczymy ponadto
$\int_0^\infty(\sin x/x)\,\mathrm dx$.

Dla $\varepsilon>0$ oznaczmy
$I(\varepsilon)=\int_0^\infty
e^{-\varepsilon x}(\sin x/x)\,\mathrm dx$.
Całka jest bezwzględnie zbieżna.
Na dowolnym przedziale
$\varepsilon\in[\delta,M]$,
$\delta>0$, pochodną podcałkową
względem $\varepsilon$
dominuje $e^{-\delta x}$.
Możemy zatem różniczkować pod całką:
\[
 I'(\varepsilon)
 =-\int_0^\infty e^{-\varepsilon x}\sin x\,\mathrm dx
 =-\frac1{1+\varepsilon^2}.
\]
Druga równość wynika z części urojonej
całki $\int_0^\infty e^{-(\varepsilon-i)x}
\,\mathrm dx=(\varepsilon-i)^{-1}$.
Ponieważ $|\sin x|/x\le1$ dla $x>0$,
mamy $|I(\varepsilon)|\le1/\varepsilon$.
Całkując równanie różniczkowe
od $\varepsilon$ do nieskończoności, otrzymujemy
$I(\varepsilon)=\pi/2-\arctan\varepsilon$.

Uzasadnijmy przejście do granicy przy
$\varepsilon\downarrow0$. Dla $0<A<B$
funkcja $g_\varepsilon(x)=e^{-\varepsilon x}/x$
jest dodatnia i malejąca do zera dla
każdego $\varepsilon\ge0$.
Całkowanie przez części daje
\[
 \left|\int_A^B g_\varepsilon(x)\sin x\,\mathrm dx\right|
 \le g_\varepsilon(A)+g_\varepsilon(B)
       +\int_A^B|g_\varepsilon'(x)|\,\mathrm dx
 =2g_\varepsilon(A)\le\frac2A.
\]
To oszacowanie jest jednostajne względem
$\varepsilon$ i pozwala kontrolować ogony.
Na każdym $[0,A]$ funkcje podcałkowe
zbiegają jednostajnie, gdy
$\varepsilon\downarrow0$, do
ciągłego przedłużenia $\sin x/x$.
Dzieląc całki na $[0,A]$ i $[A,\infty)$,
przechodzimy więc do granicy i otrzymujemy
\[
 \int_0^\infty\frac{\sin x}{x}\,\mathrm dx
 =\lim_{\varepsilon\downarrow0}I(\varepsilon)
 =\frac{\pi}{2}.
\]
Zbieżność niewłaściwej całki nie oznacza
elementarności jej funkcji pierwotnej.
\end{przyklad}

\section{Logarytm całkowy i całka wykładnicza}

\begin{twierdzenie}[Logarytm całkowy]
\label{aa:tw:logarithmic-integral-nonelementary}
Na żadnym niepustym otwartym przedziale
$I\subset(0,1)$ lub $I\subset(1,\infty)$
funkcja $1/\log x$ nie ma elementarnej
funkcji pierwotnej. Również $e^x/x$
nie ma elementarnej funkcji pierwotnej
na żadnym przedziale nieprzechodzącym
przez $0$.
\end{twierdzenie}
\begin{proof}
Najpierw rozważmy $e^x/x$.
Lemat~\ref{aa:lem:exp-polynomial-constant-field}
zapewnia, że $t=e^x$ jest
przestępny nad $\mathbb C(x)$,
$Dt=t$ i nie dodaje stałych.
Z Twierdzenia~\ref{aa:tw:hyperexponential-criterion}
wynika, że elementarna funkcja pierwotna
wymagałaby $r'+r=1/x$ dla
$r\in\mathbb C(x)$. Lemat
\ref{aa:lem:simple-pole-obstruction}
w punkcie $x=0$, przy regularnym $h=1$,
wyklucza takie rozwiązanie.
Biegun równania wymiernego jest przeszkodą
globalnej tożsamości w $\mathbb C(x)$,
nawet gdy badany przedział nie zawiera zera.

Przypuśćmy teraz, że na przedziale
$I\subset(0,1)$ lub $I\subset(1,\infty)$
istnieje elementarna funkcja pierwotna
$F$ funkcji $1/\log x$.
Podstawmy $x=e^u$ i różniczkujmy:
\[
 \frac{\mathrm d}{\mathrm du}F(e^u)
 =\frac{e^u}{u}.
\]
Skończone wyrażenie elementarne
po takim podstawieniu nadal pozostaje
elementarne; wybrane lokalne gałęzie
funkcji mają sens również w zespolonym
otoczeniu zwykłego punktu.
Dawałoby to elementarną funkcję
pierwotną $e^u/u$, którą właśnie
wykluczyliśmy. Otrzymujemy sprzeczność.
\end{proof}

\begin{przyklad}[Dwie funkcje specjalne]
\label{aa:prz:li-ei-definitions}
Na $x>1$ możemy określić
$\operatorname{Li}(x)=
\int_2^x(\log u)^{-1}\,\mathrm du$.
Odcinek całkowania nie przecina osobliwości
$x=1$; na $(0,1)$ wybiera się punkt
początkowy z tego przedziału.
W podobny sposób na $x>0$ określamy
$\operatorname{Ei}(x)$, z dokładnością
do stałej, jako funkcję pierwotną $e^x/x$.
Z twierdzenia wynika, że obie funkcje
nie mają elementarnych lokalnych wzorów.
Dla kontrastu
\[
 \int\frac1{x\log x}\,\mathrm dx
 =\log|\log x|+C
\]
na każdym z przedziałów $(0,1)$
oraz $(1,\infty)$; pochodną prawej
strony sprawdzamy przez regułę łańcuchową.
\end{przyklad}

\section{Całka eliptyczna jako odmienny rodzaj przeszkody}

Rozważmy funkcję algebraiczną $y$ daną
równaniem $y^2=1-x^4$. Nie jest ona
hiperwykładnicza nad $\mathbb C(x)$
w sensie poprzednich kryteriów:
$Dy/y=-2x^3/(1-x^4)$ wprawdzie należy
do $\mathbb C(x)$, ale $y$ jest
\emph{algebraiczny}, a nie przestępny.
Wymagana będzie odrębna metoda oparta
na formach różniczkowych i rzędach
zer oraz biegunów.

\begin{lemat}[Ciało krzywej i lokalne współrzędne]
\label{aa:lem:elliptic-charts}
Niech $K=\mathbb C(x,y)$, gdzie
$y^2=1-x^4$, $Dx=1$ i
$Dy=-2x^3/y$. Wówczas $[K:\mathbb C(x)]=2$
i $C_K=\mathbb C$. Każdy $v\in K$
ma jednoznaczną postać $A(x)+B(x)y$
z $A,B\in\mathbb C(x)$.
Na gładkiej krzywej tej relacji,
z dołączonymi dwoma punktami
nad nieskończonością, lokalne współrzędne
można wybrać następująco:
\begin{enumerate}
\item gdy $x=a$ i $y\ne0$, współrzędną
jest $z=x-a$;
\item gdy $x=a$, $y=0$ i $a^4=1$,
współrzędną jest $z=y$, a $x-a$
ma rząd $2$;
\item nad nieskończonością, w zmiennych
$s=1/x$, $w=y/x^2$,
mamy $w^2=s^4-1$ i dwa punkty
$(s,w)=(0,i)$, $(0,-i)$.
W obu współrzędną jest $s$, a
$y=w/s^2$ ma biegun rzędu $2$.
\end{enumerate}
\end{lemat}
\begin{proof}
Wielomian $1-x^4$ ma cztery różne
pierwiastki i dlatego nie jest
kwadratem funkcji wymiernej
(kwadraty mają parzyste rzędy
zer i biegunów). Stąd
$T^2-(1-x^4)$ jest nierozkładalny
w $\mathbb C(x)[T]$ i otrzymujemy
stopień rozszerzenia oraz
postać $A+By$. Ponieważ
$\mathbb C$ jest algebraicznie
domkniętym ciałem stałych
$\mathbb C(x)$, wniosek
\ref{aa:cor:algebraic-constants}
zapewnia $C_K=\mathbb C$.

Relacja algebraiczna daje
$2y\,\mathrm dy=-4x^3\,\mathrm dx$.
Gdy $y\ne0$, z twierdzenia
o funkcji uwikłanej można
lokalnie wyrazić $y$ przez $x$.
Gdy $y=0$, zachodzi $x^4=1$;
ponieważ $-4x^3\ne0$, możemy
wyrazić $x$ jako funkcję $y$.
Dokładniej, dla gałęzi mającej
$x(0)=a$ mamy
$x=a(1-y^2)^{1/4}$.
Współczynnik przy $y^2$
wynosi $-a/4\ne0$, więc
$x-a$ ma rząd dokładnie $2$.

Dla $s=1/x$ i $w=y/x^2$
po przekształceniu równania
otrzymujemy $w^2=s^4-1$.
Nad $s=0$ są dokładnie dwa
punkty $w=\pm i$. Ponieważ
$2w\ne0$, $w$ jest tam
holomorficzną, niezerową
funkcją $s$.
Wprowadzone współrzędne
obejmują wszystkie skończone
punkty krzywej oraz wszystkie
punkty nad nieskończonością.
Każdy niezerowy element $K$
ma w każdej z nich
niezerowe rozwinięcie Laurenta
o skończonym rzędzie.
\end{proof}

\begin{lemat}[Funkcje bez biegunów na krzywej]
\label{aa:lem:elliptic-no-poles}
Przy oznaczeniach
Lematu~\ref{aa:lem:elliptic-charts},
jeśli element $v\in K$
nie ma bieguna w żadnym
z opisanych tam punktów,
to $v\in\mathbb C$.
\end{lemat}
\begin{proof}
Niech $v=A(x)+B(x)y$.
Inwolucja $\tau(x,y)=(x,-y)$
przestawia punkty krzywej, także
oba punkty nad nieskończonością.
Funkcja $\tau(v)=A-By$ nie ma
zatem biegunów. To samo dotyczy
$A=(v+\tau(v))/2$ oraz
$By=(v-\tau(v))/2$.

Gdyby funkcja wymierna $A(x)$
miała biegun w pewnym skończonym
$x=a$, miałaby biegun nad $a$
na krzywej, gdyż współrzędna
$x-a$ ma dodatni rząd.
Zatem $A$ nie ma skończonych
biegunów i jest wielomianem.
Przy $s=0$ wielomian
dodatniego stopnia miałby
biegun, więc $A\in\mathbb C$.

Rozważmy $By$. W punkcie
$x=a$, gdzie $1-a^4\ne0$,
funkcja $y$ jest niezerowa
i regularna; biegun $B$
stałby się więc biegunem $By$.
W punktach $a^4=1$ mamy
$\operatorname{ord}(y)=1$
oraz $\operatorname{ord}(x-a)=2$.
Jeśli $B$ miałaby biegun
rzędu $m\ge1$ w $x=a$,
rząd $By$ wynosiłby
$1-2m<0$, znów sprzeczność.
Funkcja $B$ nie ma więc
skończonych biegunów
i jest wielomianem.
Jeżeli $B\ne0$ ma stopień $d$,
to przy $s=0$ element
$By$ ma biegun rzędu $d+2$,
ponieważ $y=w/s^2$, a
$w(0)=\pm i\ne0$.
Zatem $B=0$, a $v=A$
jest stała.
\end{proof}

\begin{lemat}[Forma $\mathrm dx/y$]
\label{aa:lem:elliptic-regular-form}
Na krzywej z
Lematu~\ref{aa:lem:elliptic-charts}
forma różniczkowa
$\omega=\mathrm dx/y$ jest
wszędzie holomorficzna
i nie jest formą zerową.
\end{lemat}
\begin{proof}
Gdy $y\ne0$ i $x$ jest
lokalną współrzędną,
$\omega=\mathrm dx/y$
jest regularna i niezerowa.
Przy $y=0$ z relacji
$2y\,\mathrm dy=-4x^3\,\mathrm dx$
otrzymujemy
$\omega=-\mathrm dy/(2x^3)$;
ponieważ $x^4=1$,
współczynnik jest regularny
i niezerowy.
Nad nieskończonością
$\mathrm dx=-s^{-2}\,\mathrm ds$
i $y=w/s^2$, zatem
$\omega=-\mathrm ds/w$.
Również ten współczynnik
jest regularny i niezerowy
przy $w=\pm i$.
\end{proof}

\begin{twierdzenie}[Nieelementarność całki eliptycznej]
\label{aa:tw:elliptic-integral-nonelementary}
Niech $K=\mathbb C(x,y)$ będzie
ciałem różniczkowym, dla którego
$y^2=1-x^4$, $Dx=1$,
$Dy=-2x^3/y$. Wtedy element
$1/y$ nie ma elementarnej
funkcji pierwotnej nad $K$
bez nowych stałych.
W szczególności funkcja
$1/\sqrt{1-x^4}$ nie ma
elementarnej funkcji pierwotnej
na żadnym niepustym otwartym
przedziale zawartym w $(-1,1)$.
\end{twierdzenie}
\begin{proof}
Załóżmy przeciwnie istnienie
elementarnej funkcji pierwotnej
w rozszerzeniu bez nowych stałych.
Na mocy Twierdzenia
\ref{aa:tw:liouville-elementary-integration}
istnieją $v\in K$,
$u_1,\ldots,u_N\in K^\times$
oraz $c_1,\ldots,c_N\in\mathbb C$,
dla których
\[
 \frac1y=Dv+\sum_{j=1}^Nc_j\frac{Du_j}{u_j}.
\]
Pomnożenie przez $\mathrm dx$
daje tożsamość form na krzywej
\begin{equation}\label{aa:eq:elliptic-liouville-form}
 \omega=\mathrm dv+
       \sum_{j=1}^Nc_j\frac{\mathrm du_j}{u_j}.
\end{equation}
Jest to tożsamość elementów
ciała funkcji krzywej
i ich różniczek, więc
możemy badać ją w każdym
lokalnym układzie współrzędnych.

\emph{Eliminacja części dokładnej.}
W lokalnej współrzędnej $z$
element $u\in K^\times$ ma
postać $u=z^m g(z)$,
gdzie $m\in\mathbb Z$,
$g$ jest holomorficzna
i $g(0)\ne0$.
Wtedy
\[
 \frac{\mathrm du}{u}
 =m\frac{\mathrm dz}{z}
   +\frac{\mathrm dg}{g}.
\]
Pochodne logarytmiczne
mają zatem co najwyżej
bieguny proste.
Jeśli $v$ miałby
biegun rzędu $q\ge1$,
to $\mathrm dv$
miałaby biegun rzędu
$q+1\ge2$:
wyraz wiodący jest
$-q a_{-q}z^{-q-1}\,\mathrm dz\ne0$.
Nie mogłyby usunąć go
pochodne logarytmiczne,
a $\omega$ nie ma biegunów
na mocy
Lematu~\ref{aa:lem:elliptic-regular-form}.
Zatem $v$ nie ma biegunów
w żadnym punkcie. Lemat
\ref{aa:lem:elliptic-no-poles}
daje $v\in\mathbb C$,
czyli $\mathrm dv=0$.

\emph{Eliminacja części logarytmicznej.}
Współczynnik przy
$\mathrm dz/z$, czyli residuum
formy $\mathrm du_j/u_j$
w punkcie $P$, wynosi
$\operatorname{ord}_P(u_j)\in\mathbb Z$.
Ponieważ $\omega$ nie ma
biegunów, z
\eqref{aa:eq:elliptic-liouville-form}
wynika dla każdego punktu $P$
\[
 \sum_{j=1}^Nc_j\operatorname{ord}_P(u_j)=0.
\]
Jeśli $N=0$ albo wszystkie
$c_j=0$, prawa strona
\eqref{aa:eq:elliptic-liouville-form}
już jest zerowa.
W przeciwnym razie
wybierzmy bazę
$\gamma_1,\ldots,\gamma_s$
nad $\mathbb Q$
przestrzeni liniowej
rozpiętej przez
$c_1,\ldots,c_N$.
Zapiszmy
$c_j=\sum_{\nu=1}^s
q_{\nu j}\gamma_\nu$,
$q_{\nu j}\in\mathbb Q$.
Niezależność $\gamma_\nu$
nad $\mathbb Q$ daje
dla wszystkich $P$ i $\nu$
równości
$\sum_jq_{\nu j}\operatorname{ord}_P(u_j)=0$.
Niech $M\ge1$ będzie wspólnym
mianownikiem wszystkich
$q_{\nu j}$.
Dla każdego $\nu$ element
$U_\nu=\prod_{j=1}^N
u_j^{Mq_{\nu j}}\in K^\times$
ma całkowite wykładniki i
$\operatorname{ord}_P(U_\nu)=0$
dla każdego punktu $P$.
Nie ma zatem biegunów;
z Lematu~\ref{aa:lem:elliptic-no-poles}
mamy $U_\nu\in\mathbb C^\times$,
więc $\mathrm dU_\nu=0$.
Wobec tego
\[
 \sum_{j=1}^Nc_j\frac{\mathrm du_j}{u_j}
 =\frac1M\sum_{\nu=1}^s
       \gamma_\nu\frac{\mathrm dU_\nu}{U_\nu}
 =0.
\]
Równość~\eqref{aa:eq:elliptic-liouville-form}
dawałaby $\omega=0$,
wbrew Lematowi
\ref{aa:lem:elliptic-regular-form}.
Otrzymujemy sprzeczność.

Na przedziale $I\subset(-1,1)$
wybór dodatniej gałęzi
$y=\sqrt{1-x^4}$ daje
zanurzenie $K$ w ciele
zarodków funkcji meromorficznych.
Gdyby istniał tam lokalny
elementarny wzór na funkcję
pierwotną, dawałby
elementarne rozszerzenie
różniczkowe bez nowych
stałych nad $K$.
Właśnie wykluczyliśmy
takie rozszerzenie.
\end{proof}

\begin{uwaga}[Rola residuów]
\label{aa:uw:elliptic-residue-warning}
Nie wystarczyłoby powiedzieć,
że holomorficzna forma
ma zerowe residua:
niezerowa forma
$\omega=\mathrm dx/y$
ma właśnie tę własność.
Kluczowy krok dowodu
wykorzystuje \emph{całkowitość}
rzędów zer i biegunów funkcji
$u_j$ oraz niezależność
współczynników $c_j$
nad $\mathbb Q$.
Pozwala to skonstruować
rzeczywiste elementy $U_\nu\in K$,
które są stałe, i w konsekwencji
usunąć całą sumę
pochodnych logarytmicznych.
\end{uwaga}

\section{Przykłady i zakres zastosowań}

\begin{przyklad}[Kontrast z funkcjami wymiernymi]
\label{aa:prz:rational-primitive}
Każda $f\in\mathbb C(x)$
ma lokalną elementarną
funkcję pierwotną.
Istotnie, rozkład
na ułamki proste ma postać
\[
 f=P(x)+\sum_a\sum_{k=1}^{m_a}
     \frac{c_{a,k}}{(x-a)^k},
 \qquad P\in\mathbb C[x].
\]
Wielomian i wyrazy z $k\ge2$
całkujemy wymiernie, korzystając
z $D((x-a)^{1-k})=(1-k)(x-a)^{-k}$.
Dla $k=1$ funkcją pierwotną
jest $c_{a,1}\log(x-a)$
na ustalonej lokalnej gałęzi.
Ta obserwacja uzasadniała
brak warunku dla $s_0$
w Twierdzeniu
\ref{aa:tw:laurent-hyperexp-criterion}.
\end{przyklad}

Zebrane przykłady pokazują
dwie odrębne przeszkody
w elementarnym całkowaniu:
równanie wymierne
$r'+hr=s$ może nie mieć
rozwiązania, nawet gdy
całkowana funkcja jest
całkowita lub ma prosty
opis wykładniczy,
natomiast całka eliptyczna
wymaga analizy form
różniczkowych na
krzywej algebraicznej.
W następnym rozdziale,
dopiero po stabilizacji
niniejszego, wykorzystamy
ciała różniczkowe do
badania automorfizmów
rozszerzeń Picarda--Vessiota
równań liniowych.
